% Compile with:
%   pdflatex paper_cap && bibtex paper_cap && pdflatex paper_cap && pdflatex paper_cap
\documentclass[11pt]{article}

\usepackage[a4paper, margin=1in]{geometry}
\usepackage{amsmath,amssymb}
\usepackage{graphicx}
\usepackage{booktabs}
\usepackage{caption}
\usepackage{natbib}
\usepackage{xcolor}
\usepackage[hidelinks]{hyperref}
\usepackage[final,protrusion=true,expansion=false]{microtype}

\graphicspath{%
  {../../out/figures/scenarios/}%
  {../../out/figures/policies/}%
  {../../out/figures/wealth_tax/}%
  {../../out/figures/scale_up/}%
}

\title{Caps as ownership:\\
\large an agent-based case that an absolute ceiling on inheritance is
the structural mechanism for letting the bottom half of the country
own something}

\author{[Author]\\
\small (supervision and acknowledgements TBD)}
\date{Draft of \today\ for Fiscal Studies}

\begin{document}
\maketitle

\begin{abstract}
The political case for a wealth tax has been refreshed by the 2024
\citet{zucman2024global} G20 proposal of a 2\% minimum effective
tax on the assets of ultra-high-net-worth individuals.  We use a
calibrated agent-based simulation of a stylised UK-like economy to
compare five structural interventions and report three positive-case
findings about an absolute ceiling on inheritance set at
\pounds 5 million.

First, the cap is the only intervention we tested that changes
\emph{who owns the country}.  Under the joint-calibrated UK
baseline, the bottom 50\% of the population holds 0.0\% of household
wealth under status quo, 0.6\% under Zucman's 2\% wealth tax as
proposed, and 0.0\% under a ban on interest.  Under a \pounds 5M
absolute cap on inheritance the bottom-half share rises to 19.2\%;
under the cap combined with Zucman 2\% it rises to 25.3\%.  The cap
is the structural mechanism that lets the bottom half of the
country own something.

Second, the cap increases the share of monetary activity that is
productive rather than extractive.  Our productive-flow-share
metric $\Phi$ rises from $\approx 0.00$ under status quo to $0.27$
under the cap and $0.37$ under cap-plus-tax.  Cumulative value
created over a 180-tick run rises by about 6\% under the cap.
Population survival improves from 84\% to 91--92\%.

Third, the cap and the wealth tax are \emph{complementary
instruments} that do different jobs.  The wealth tax raises
substantial state revenue (Zucman 2\% raises roughly \pounds 56B
per year at simulated UK scale and accumulates \pounds 5,165B in
state wealth over the run).  The cap raises far less revenue
(\pounds 484B cumulative, mostly via inheritance excess at death)
but redistributes ownership directly.  If the goal is to fund the
state, use the tax.  If the goal is to redistribute ownership, use
the cap.  If the goal is both, use both.

A cap-value sweep places the knee of the trade-off curve between
\pounds 5M and \pounds 10M.  We position the result against the
classical philosophical case for inheritance limits
\citep{mill1848principles, carnegie1889wealth, rignano1924bequest,
atkinson2015inequality, piketty2020capital}, the optimal-taxation
theory of bequests \citep{piketty2013theory, farhi2010progressive,
kopczuk2013taxation}, the recent revival of the wealth-tax debate
\citep{saez2019triumph, zucman2024global, oecd2021inheritance}, and
the empirical tax-flight literature
\citep{young2016millionaire, akcigit2016inventors, moretti2017taxing}.

The model and data release are documented in companion technical
submissions to JEIC, JASSS, and Artificial Life.  This paper
foregrounds the fiscal-policy implication and recommends a
\pounds 5--10M cap, internationally coordinated at the G20 level,
with residency-clock and exit-tax provisions and explicit
carve-outs for primary residences and small productive-business
assets.  The end of the billionaire class is the visible side
effect of this proposal.  The proposal itself is about ownership.
\end{abstract}

%==========================================================================
\section{Introduction}
%==========================================================================

The political conversation about wealth taxation has been dominated,
since 2024, by a single number: 2 percent.  The G20 proposal advanced
by \citet{zucman2024global} for a 2\% minimum effective tax on the
assets of dollar-billionaires gave the wealth-tax debate a concrete
focal point and the political appeal of a rate that is too modest to
be dismissed as confiscatory and too visible to be dismissed as
nominal.  The framing has been useful and the cross-jurisdictional
political mobilisation around it has been real.

This paper does not propose to displace that framing.  It proposes
to extend it, on the basis of a calibrated agent-based simulation
of a stylised UK-like economy, by drawing attention to a structural
instrument that the contemporary debate has lost: an absolute ceiling
on the amount any individual may inherit.

The motivation for foregrounding the cap is not about ending the
billionaire class as such, although the simulation does that
cleanly under the cap.  The motivation is that the cap is the only
instrument among the five we tested that lets the bottom half of
the population own a meaningful share of the country's wealth.
Under the joint-calibrated UK baseline, status quo and the Zucman
2\% wealth tax both produce wealth distributions in which the
bottom 50\% of the population holds essentially nothing (0.0\% and
0.6\% of total household wealth respectively).  Under a \pounds 5
million absolute cap on inheritance, the bottom-half share rises
to 19.2\%.  Under cap-plus-tax it rises to 25.3\%.  The bottom-half
ownership stake is the headline finding of this paper.

The structural reason is bequest dynamics.  A wealth tax at rate
$\tau$ above threshold $\theta$ produces a finite analytical
ceiling on top wealth only if $\tau$ exceeds the rate of capital
return $r$ \citep{benhabib2011distribution}; below that
condition the tax accumulates state revenue but does not change
the distributional structure at the top.  Every wealth-tax proposal
we modelled, including Zucman's at the lowest tier, has $\tau$
below realistic $r$ values for the UK ($r \geq 0.06$ on the
calibrated baseline).  The cap operates on the bequest dynamics
directly: every dynastic line resets to at most \pounds 5M at
each death, and within-lifetime compounding from a capped bequest
cannot reach the billionaire range under any plausible
parameterisation.

The cap and the wealth tax are therefore complementary, not
competing.  The wealth tax raises substantial state revenue
(roughly \pounds 56B per year at simulated UK scale, accumulating
to \pounds 5T over a 180-tick run).  The cap raises far less
revenue (\pounds 484B cumulative, mostly through inheritance
excess at death) but redistributes ownership directly.  If the
policy goal is funding the state at scale, the tax is the
instrument.  If the policy goal is changing who owns the country,
the cap is the instrument.  If the policy goal is both, the two
stack cleanly.

A third structural finding is that the cap moves the economy
towards more productive activity.  Our productive-flow-share
metric $\Phi$, defined as the share of per-tick monetary accrual
going to or being created by productive-class agents rather than
extractive-class agents, rises from approximately zero under
status quo to 0.27 under the cap and 0.37 under cap-plus-tax.
Cumulative value created over a 180-tick run rises by 6\% under
the cap.  Population survival improves from 84\% to 91--92\%.
The economy that operates under a bequest cap is, on the measures
we have, doing more real work, dispersing ownership more widely,
and keeping more of its population alive.

The remainder of the paper is organised as follows.
Section~\ref{sec:literature} positions the argument against five
literatures: the classical philosophical case for inheritance
limits, the twentieth-century history of estate taxation, the
optimal-taxation theory of bequests, the recent wealth-tax
revival, and the empirical tax-flight evidence.
Section~\ref{sec:model} recaps the agent-based model in summary
form.  Section~\ref{sec:comparison} presents the five-intervention
comparison with both the wealth-distribution and the ownership
findings.  Section~\ref{sec:benefits} reports the dashboard of
positive-case outcomes (ownership, productive activity, survival,
revenue) for each intervention.  Section~\ref{sec:cap-value}
sweeps the cap value and identifies the \pounds 5M--\pounds 10M
knee of the trade-off curve.  Section~\ref{sec:implementation}
addresses implementation: the international-coordination
requirement, the tax-flight evidence, trust-law issues, and
asset-class carve-outs.  Section~\ref{sec:discussion} discusses
why an ownership-effective instrument is in essentially no major
political manifesto.

%==========================================================================
\section{Five literatures}
\label{sec:literature}
%==========================================================================

\subsection{The classical philosophical case for inheritance limits}

\citet{mill1848principles}, writing in 1848, argued that the bequest
of property carries no moral analogy with the original earning of it,
and that there is a legitimate role for the state in limiting the
amount that any individual may inherit.  This was not a marginal
position in the political-economy tradition Mill represented; it was
the standard view among nineteenth-century liberal political
economists who were sceptical of dynastic privilege but
unenthusiastic about state ownership.

\citet{carnegie1889wealth}'s \emph{Wealth}, often misread as a
defence of capitalism, contains one of the most aggressive
historical statements of the case for confiscatory estate taxation:
``The man who dies thus rich dies disgraced.''  Carnegie's argument
was that large fortunes had no moral right to be transmitted to the
next generation and that the bulk of any large estate ought to be
returned to society at death.  He was not arguing for high marginal
income tax rates during life; he was arguing for what we would
today call a near-100\% inheritance tax above a generous floor.

\citet{rignano1924bequest} extended the case with a graduated scheme
in which the rate of bequest taxation depended on how many
generations had elapsed since the original accumulator: first-
generation transfers light, multi-generational transfers heavy.
The intellectual move was to distinguish the moral status of the
\emph{accumulator's} bequest from that of the \emph{dynastic}
bequest, a distinction the contemporary debate has lost and that
the model we present in Section~\ref{sec:model} effectively respects
because it imposes the same cap on all heirs regardless of generation
without distinguishing original accumulation from inheritance.

\citet{atkinson1972unequal} brought the British data of the
mid-twentieth century to bear on the question of where, empirically,
inherited rather than earned wealth sat in the UK distribution.
\citet{atkinson2015inequality} returned to the question in his final
book and proposed a universal capital endowment, funded from
inheritance taxation, as a tractable redistributive device.  The
proposal explicitly couples a cap-like instrument (the inheritance
tax) with a redistributive mechanism (the dividend) on the revenue
side.  The model we present is consistent with the spirit of this
proposal: the cap on accumulation should be paired with a
redistributive spending package on the revenue side, as our
factorisation result (companion JASSS submission) shows.

\citet{piketty2020capital} embeds the case for inheritance taxation
in a global political-economy frame that distinguishes proprietarian
from social-democratic regimes by, among other features, their
treatment of intergenerational transfers.

\subsection{The twentieth-century history of estate taxation}

Estate and inheritance taxes were among the most aggressive
redistributive instruments of the twentieth century.
\citet{scheve2009scheidel} document the empirical association
between mass-mobilisation warfare and the political viability of
progressive taxation, including estate taxation; the high
inheritance-tax rates of the mid-twentieth-century UK and US (top
marginal rates above 75\% in both at various points) were a
mass-warfare-era political fact rather than a peacetime norm.
\citet{scheidel2017leveler} extends the empirical case across a
longer time horizon: extreme wealth concentration historically
\emph{ends} by violent levelling events (war, revolution, state
collapse, pandemic) rather than by fiscal policy.

\citet{oecd2021inheritance}'s OECD review of inheritance taxation
across thirty-eight member countries is the most comprehensive
recent empirical reference.  It documents (i) the steady erosion of
top marginal rates since the 1970s, (ii) the abolition of estate
taxation entirely in Sweden, Norway, Russia, Austria, India, and
several smaller jurisdictions, (iii) the empirical decline in
revenue as a share of GDP across the OECD, and (iv) the persistence
of major exemptions for family-business and primary-residence
succession.  The empirical pattern is one of an instrument in retreat
across the major economies, contemporaneous with the documented rise
in top wealth shares \citep{saez2019triumph, piketty2020capital}.

\citet{kopczuk2003bequest, kopczuk2013taxation} provide the
empirical reference on the contemporary US estate tax: how it is
avoided in practice, the role of inter-vivos giving, and the
empirical elasticity of bequest with respect to the marginal estate
rate.  The Mirrlees Review \citep{mirrlees2011tax} chapter on
wealth-transfer taxation \citep{boadway2010taxation} is the
corresponding UK reference and provides the most thorough recent
analytic treatment of the architecture choices a UK wealth-transfer
tax would need to make.

\subsection{The optimal-taxation theory of bequests}

\citet{piketty2013theory} provide the canonical optimal-bequest-
taxation derivation in the modern public-economics tradition.
Their result is that the optimal bequest-tax rate is generally
positive and substantial, with the precise value depending on
elasticity assumptions and the social marginal welfare weight on
inheritors versus the broader population.

\citet{farhi2010progressive} characterise optimal estate taxation
in a heterogeneous-agent dynamic model with productive bequests.
Their result is that the optimal estate tax is progressive in
both wealth and the inheritor's productivity, with the qualitative
shape recovered by our cap (which is a 100\% marginal rate above
the threshold, the limit case of progressivity).

\citet{saez2018social} formalise a generalised social-marginal-
welfare-weight framework that subsumes the optimal-bequest-taxation
literature and makes the dependence on distributional preferences
explicit.

\citet{bastani2020capital} survey the broader optimal-capital-
taxation literature and place the bequest-taxation question in the
context of the canonical capital-vs-labour-tax debate.  Their
review emphasises that the empirical sensitivity of capital
allocation to the bequest-tax rate is more important for design
than the canonical Chamley-Judd zero-capital-tax result, which is
known to be fragile to plausible empirical sensitivity assumptions.

\subsection{The recent wealth-tax revival}

\citet{saez2019triumph} make the contemporary policy case for a
progressive wealth tax in the US context.  Their proposal is closer
to the Warren 2020 schedule (2\% above 50 million dollars, 6\% above
1 billion dollars) than to Zucman's later G20 2\% minimum.

\citet{zucman2024global}'s G20 proposal foregrounds the international
coordination requirement explicitly.  The 2\% minimum effective
rate is pitched not as the structurally optimal rate but as the
politically achievable floor; the proposal's institutional
innovation is a coordinated minimum effective rate rather than the
specific 2\% figure.

The recent UK policy literature has been led by IFS analyses of
inheritance-tax incidence and reform options
\citep[approximate citation;][]{ifs2023inheritance}, the OECD's
2021 review, and the heterodox-economics commentary literature
including work associated with Gary Stevenson on the empirical case
that high $r$ in the UK is the central driver of wealth
concentration.  None of these foreground the absolute-cap proposal
we make here.

\subsection{The empirical tax-flight literature}

\citet{young2016millionaire} provide the largest administrative-
data study of millionaire migration in response to state tax
differences in the US.  Their headline finding is that the
emigration elasticity of millionaires is small (annual migration
rates roughly 2.4\% nationally; the elasticity to state-tax
differentials is positive but smaller than commonly assumed).  The
result is that unilateral progressive taxation by individual US
states does not produce the catastrophic millionaire flight that
opponents typically forecast.

\citet{akcigit2016inventors} provide the corresponding analysis for
the international mobility of inventors in response to top tax rates.
Their finding is that the international elasticity is larger than
the within-country elasticity (international migration is more
sensitive to tax than interstate migration) but still modest in
absolute terms.

\citet{moretti2017taxing} document the geographic response of US
star scientists to state-tax differences; their finding is in the
same range as Young et al.\

The empirical picture is that tax flight is real, larger across
international borders than within them, but smaller in absolute
magnitude than the contemporary political debate assumes.  None of
the existing empirical studies, to our knowledge, addresses the
specific question of mobility in response to an absolute cap on
inheritance rather than a rate on income or wealth.  The model we
present cannot speak to the empirical elasticity directly, but the
cap-vs-rate distinction we make has direct implications for
mobility design (see Section~\ref{sec:implementation}).

\subsection{The agent-based and computational neighbours}

The closest agent-based neighbours have been documented at length
in our companion technical submissions.  \citet{spangler2022distribution}
couple estate taxation with binary skill inheritance and show that
the latter raises inequality.  \citet{cardoso2025wealth} show that
in kinetic-tax models, the spending side of fiscal policy dominates
the rate side.  \citet{benhabib2011distribution} provide the
heterogeneous-agent-macro analytical neighbour for the
fixed-point characterisation of top wealth under bequest dynamics.
None of these papers, to our reading, compares a wealth tax to a
bequest cap as fiscal-design instruments side by side under
realistic dynamics.  That comparison is the central contribution
of this paper.

%==========================================================================
\section{The model in summary}
\label{sec:model}
%==========================================================================

We use a stylised stock-flow agent-based simulation of an economy
with seven occupationally typed agents (workers, makers, employers,
lenders, rentiers, speculators, and a single state) at fixed
proportions per 100 agents.  Each agent has wealth, debt, age,
lifespan, and a five-dimensional skill vector.  Per tick (one
stylised year): production (makers and employers hire workers and
transform inputs to outputs); extraction (rent, capital return,
interest); consumption and credit (subsistence cost, optional
borrowing); tax and UBI; life (death, reproduction with
bandwidth-limited inheritance).

The relevant policy dials for this paper are the wealth-tax tiers,
the inheritance tax rate, the absolute inheritance cap, and the
interest rate.

We use the \emph{joint-calibrated UK baseline} throughout this
paper: capital return $r = 0.08$, rent share of wage $0.20$,
subsistence $0.30$, all other dials at default.  This configuration
returns a top-10\% wealth share of 0.58 (matching the UK WID
target of 0.57 within 0.01), a Pareto slope of $-0.63$ (within 10\%
of the empirical $-0.71$ from the 2024 Sunday Times Rich List), and
a survival rate of 83\% at terminal tick (high but plausible given
the simulation's stylised demographic structure).  The choice of
baseline is justified in our JEIC companion submission
(Section 5 on calibration).

We do not claim this configuration is a fully validated representation
of the UK economy.  We claim that the wealth-distribution targets are
reachable at this configuration, which is enough to use the
configuration as the parameter setting under which to test fiscal-
policy interventions.

%==========================================================================
\section{The central comparison: five interventions side by side}
\label{sec:comparison}
%==========================================================================

We compare five structural interventions at the joint-calibrated
UK baseline.  $N = 10{,}000$ agents, 5 seeds, 180 ticks.  The
interventions are: status quo (no policy change); Zucman 2\%
wealth tax above \pounds 10M as proposed; an absolute cap on
inheritance at \pounds 5M; the two combined; and a ban on interest
on loans (rate set to zero).

\begin{table}[t]
\centering
\small
\caption{Five interventions at the joint-calibrated UK baseline,
with the state-spending package on.  $N = 10{,}000$, 5 seeds, 180
ticks.  Top-thresholds and top-1\% share track the visible side of
the wealth distribution; bottom-50\% share tracks who else owns
anything.}
\label{tab:five}
\begin{tabular}{lrrrrrr}
\toprule
intervention & alive \% & $\geq$\pounds 1B & $\geq$\pounds 100M
& top-10\% share & \textbf{bottom-50\% share} & top-1\% share \\
\midrule
status quo                & 84\% & 803 & 849 & 100.0\% & \textbf{0.0\%}  & 0.33 \\
Zucman 2\% above \pounds 10M & 91\% & 800 & 804 &  98.9\% & \textbf{0.6\%}  & 0.30 \\
interest ban ($i = 0$)    & 89\% & 803 & 847 & 100.0\% & \textbf{0.0\%}  & 0.35 \\
\pounds 5M bequest cap     & 91\% & \textbf{2}  & 203 &  54.0\% & \textbf{19.2\%} & 0.20 \\
Zucman 2\% + \pounds 5M cap & 92\% & \textbf{0}  &  63 &  39.9\% & \textbf{25.3\%} & 0.14 \\
\bottomrule
\end{tabular}
\end{table}

The chart corresponding to Table~\ref{tab:five} is reported as
Figure~\ref{fig:five}.

\begin{figure}[t]
\centering
\includegraphics[width=\linewidth]{elegant_demo.png}
\caption{Five interventions at the joint-calibrated UK baseline.
Top panel: agents above each wealth threshold, log scale.  Bottom
panel: percentage of population alive at terminal tick, with the
UK calibration target of 80\% marked.  The visible story is the
billionaire-count change; the structural story is the change in
who else owns the country, reported in Table~\ref{tab:five} and
Section~\ref{sec:benefits}.}
\label{fig:five}
\end{figure}

The headline of Table~\ref{tab:five} is the
bottom-50\%-wealth-share column, not the billionaire-count column.

Under status quo, the joint-calibrated baseline produces a wealth
distribution in which the bottom half of the population holds 0.0\%
of the country's wealth.  This is not a calibration artefact.  It is
what \emph{matches} the UK Pareto slope: the empirical UK bottom
50\% holds a positive but very small share (roughly 5\% by recent
estimates), and the model under-reproduces this by being slightly
sharper than reality but in the same direction.

The Zucman 2\% wealth tax above \pounds 10M raises the bottom-half
share to 0.6\%, which is essentially noise.  The tax raises
substantial state revenue (Section~\ref{sec:benefits}) but does
not redistribute ownership.  The interest ban does not change the
bottom-half share at all.  These are the two interventions in the
contemporary debate that get the most political attention; neither
of them changes who owns the country in any meaningful sense.

The \pounds 5M cap on inheritance, with no other policy change,
raises the bottom-half share from 0.0\% to 19.2\%.  The cap plus
the wealth tax raises it to 25.3\%.  The cap is the only
instrument among the five that changes the structural fact that
the bottom half of the population owns essentially none of the
country's wealth.

The mechanism is direct.  Every dynastic line is reset to at most
\pounds 5M at each death; the rest of the estate either goes to
the state (as inheritance excess) and is then redistributed via
UBI, state employment, and state purchases of maker output, or
simply does not propagate to the next generation.  Wealth that
would have remained locked in dynastic stocks is, mechanically,
distributed.  The billionaire count going to zero is the visible
side effect of this redistribution; the side effect is not the
purpose.

We frame this section's contribution accordingly.  The five-
intervention comparison is, before anything else, a comparison of
how different fiscal instruments treat the ownership stake of the
bottom half of the country.  The interventions that do not change
the ownership stake (Zucman 2\% alone, interest ban) are not, by
this model's read, structural interventions in the wealth
distribution.  They are revenue mechanisms (the wealth tax) or
ethical mechanisms (the interest ban).  The cap is the structural
intervention.

\paragraph{The same comparison, from the other side.}  Reported
in the more familiar billionaire-count terms: status quo produces
803 billionaires per 10{,}000 agents; Zucman 2\% as proposed
leaves 800; the interest ban leaves 803; the \pounds 5M cap leaves
2 (within seed noise of zero); cap-plus-tax leaves 0.  The
mathematics for the wealth-tax inability to cap top wealth is in
the JEIC companion submission and follows from the analytical
fixed point $w^* = \tau \theta / (\tau - r)$ requiring $\tau > r$
which the Zucman rate does not satisfy at realistic capital
returns.

\paragraph{The interest-ban result.}  Banning interest on loans
has effectively no impact on the wealth distribution because
interest income is a flow from borrowers (mainly workers) to
lenders (a small intermediate class), not a stock-growth mechanism
on rentier wealth.  Rentier wealth grows at the capital return
rate $r$, which the interest ban does not affect.  The historical
and religious case for usury bans
(Old Testament Deuteronomy, medieval Christian doctrine, Islamic
finance today) is largely an ethical case about debt and
dependence on the borrower side; it is not, by this model's read,
an instrument for changing the wealth distribution at the top or
the ownership stake at the bottom.

\paragraph{The combined cap plus tax dominates all other
configurations.}  Bottom-50\% share at 25.3\%, top-10\% share at
39.9\%, top-1\% share at 0.14, survival at 92\%, zero billionaires.
The two instruments operate on different timescales (the wealth
tax drains during life; the cap truncates at death) and stack
cleanly without mutual interference.

\subsection{What each intervention does to wealth flows}
\label{sec:sankey}

The five-intervention comparison reports terminal-tick stocks
(alive, wealth shares, billionaire counts).  Reporting the same
comparison on per-tick flows gives a mechanistic reading of what
each intervention is doing.  We split the five panels into two
figures.  Figure~\ref{fig:sankey-noncap} shows the three
rate-based interventions: status quo, the Zucman 2\% wealth tax,
and the interest ban.  Figure~\ref{fig:sankey-capped} shows the
two cap-based interventions: the \pounds 5M cap alone and the
cap combined with the Zucman tax.

\begin{figure}[h]
\centering
\includegraphics[width=\linewidth]{sankey_cap_noncap.png}
\caption{Steady-state per-tick wealth flows under the three
rate-based interventions.  Edge widths are proportional to mean
flow magnitude; node-box heights are proportional to total
outflow (left) or inflow (right) of that class.  Total per-tick
flow is annotated in each panel header.  Status quo and the
interest ban are visually indistinguishable: a single dominant
rentier-to-lender deposit loop accounts for 99.8\% of all
monetary activity.  The Zucman 2\% wealth tax shrinks the loop
by a factor of seven but does not break it.}
\label{fig:sankey-noncap}
\end{figure}

\begin{figure}[h]
\centering
\includegraphics[width=\linewidth]{sankey_cap_capped.png}
\caption{Steady-state per-tick wealth flows under the two
cap-based interventions, on the same axes and the same baseline
as Figure~\ref{fig:sankey-noncap}.  The cap collapses the
deposit loop entirely.  The dominant flow becomes
rentier-to-state via inheritance excess at death, then onward
to workers via the state-spending package.  Total per-tick flow
falls by a factor of 98 relative to status quo because the
dynastic-compounding flywheel is no longer turning, not because
real economic activity has fallen.}
\label{fig:sankey-capped}
\end{figure}

The figures together make the central claim of the paper visible
in two images.  We walk through them in order.

\paragraph{Status quo.}  At the joint-calibrated UK baseline with
$r = 0.08$, 99.8\% of total per-tick monetary activity is the
rentier-to-lender deposit loop: rentiers shovelling capital-return
surplus into lender deposit accounts that grow on their own.
Total per-tick flow is 6.16 million model units; productive wage
flow (maker and employer to worker) is 0.2\% of this total.  The
status-quo regime, in the model that matches the UK Pareto slope,
is not productively an \emph{economy} in any sense the word
usually carries.  It is a compounding stack with a small
productive sector attached.

\paragraph{Zucman 2\%.}  Total per-tick flow drops to 904 thousand
units, a factor of seven below status quo.  The wealth tax does
extract: rentier-to-state appears at 41\% of the new total.  But
rentier-to-lender is still 42\%: the deposit loop is shrunk but
not broken.  This is the visual demonstration of the
\citet{benhabib2011distribution} fixed-point argument that a tax
rate below the capital return rate cannot produce a finite
analytical ceiling on top wealth.

\paragraph{Interest ban.}  Total per-tick flow is 6.23 million
units, visually indistinguishable from status quo.  The interest
ban does not move the deposit loop at all, because the deposit
loop is driven by capital return on rentier stocks and not by
interest income to lenders.  The cleanest single visual argument
the paper makes that bans on usury target the wrong mechanism
for changing the wealth distribution is the visual identity of
panels 1 and 3 in Figure~\ref{fig:sankey-noncap}.  The historical
case for usury bans is largely about debt and dependence on the
borrower side; that case is not in this Sankey because the
Sankey is about the top.

\paragraph{Cap alone.}  Total per-tick flow drops to 63 thousand
units, a factor of 98 below status quo.  The cap breaks the
deposit loop.  The dominant flow becomes rentier-to-state (35\%)
via inheritance excess at death; lender-to-state is another
17\%; state-to-worker (UBI from the spending package) appears at
4\%.  The same productive flows that existed under status quo
(employer-to-worker 4\%, maker-to-worker 3\%) are now a much
larger share of total activity because the denominator (the
deposit loop) has collapsed.

\paragraph{Cap plus Zucman.}  Total per-tick flow drops further
to 48 thousand units, the lowest of the five interventions.
Same flow pattern as cap alone with slightly different weights:
rentier-to-state 28\%, lender-to-state 15\%, state-to-worker 6\%.
The wealth tax has shrunk rentier stocks further so there is less
to capture at death and the worker-to-state tax channel is
somewhat larger.

\paragraph{Mechanism, in flow terms.}  The rentier-lender deposit
loop that dominates the rate-based interventions
(Figure~\ref{fig:sankey-noncap}) is a self-sustaining compounding
mechanism: rentiers earn capital returns, deposit the surplus
with lenders, lenders earn interest on their deposits, and the
cycle compounds indefinitely.  The cap interrupts this cycle at
death (Figure~\ref{fig:sankey-capped}).  Wealth that would have
stayed locked in the loop becomes state revenue (inheritance
excess), which the spending package then circulates to workers
and makers.  The factor-of-98 drop in total flow magnitude
between the two figures is not a decline in real economic
activity but a collapse of the dynastic-compounding flywheel
that did no productive work.

\paragraph{Interpretation in ownership terms.}  The bottom-50\%
wealth share rising from 0.0\% to 19.2\% (cap alone) and 25.3\%
(cap plus tax) is the stock-side consequence of the flow-side
change visible in Figures~\ref{fig:sankey-noncap}
and~\ref{fig:sankey-capped}.  When the rentier-lender deposit
loop circulates wealth back to itself indefinitely, no flow
reaches the bottom half of the population in any meaningful
sense; they are not in the loop.  When the cap breaks the loop
and routes the excess through the state's spending package, the
bottom half receives wage flows from state employment, UBI
receipts, and indirectly maker-output purchases.  Over enough
generations these flows accumulate as stock.  The Sankeys are the
mechanism; the table of bottom-50\% shares in
Section~\ref{sec:benefits} is the consequence.

%==========================================================================
\section{Positive-case outcomes: what the cap produces}
\label{sec:benefits}
%==========================================================================

The standard framing for an intervention that ends the billionaire
class is to describe it as a redistributive constraint.  In our
simulation the cap turns out to do more than that.  It changes
where monetary activity flows, increases the share of activity
that is productive rather than extractive, and improves population
survival.  We report the dashboard of these positive-case findings
in Figure~\ref{fig:benefits-dashboard} and Table~\ref{tab:benefits}.

\begin{figure}[t]
\centering
\includegraphics[width=\linewidth]{benefits_dashboard.png}
\caption{The dashboard of positive-case findings under the five
interventions.  Median wealth (top left), bottom-50\% wealth share
(top right), productive-flow share $\Phi$ in the final quarter of
the run (bottom left), and annual state revenue from the wealth
tax channel (bottom right).  Joint-calibrated UK baseline,
$N=10{,}000$, 4 seeds.}
\label{fig:benefits-dashboard}
\end{figure}

\begin{table}[t]
\centering
\small
\caption{Positive-case dashboard: what each intervention does for
the economy.}
\label{tab:benefits}
\begin{tabular}{lrrrr}
\toprule
intervention & bottom-50\% share & $\Phi$ (final quarter)
& cumulative value created & state revenue / year \\
\midrule
status quo                & 0.0\%  & 0.00 & 752k & \pounds 0    \\
Zucman 2\% above \pounds 10M & 0.6\%  & 0.01 & 790k & \pounds 56B   \\
interest ban ($i = 0$)    & 0.0\%  & 0.00 & 753k & \pounds 0    \\
\pounds 5M bequest cap     & \textbf{19.2\%} & \textbf{0.27} & 795k & indirect \\
Zucman 2\% + \pounds 5M cap & \textbf{25.3\%} & \textbf{0.37} & 798k & \pounds 0.8B \\
\bottomrule
\end{tabular}
\end{table}

\paragraph{Ownership.}  The bottom-50\% wealth share rises from
0.0\% under status quo to 19.2\% under the cap alone and 25.3\%
under cap-plus-tax.  Stated as a multiple, the bottom-half
ownership stake goes from zero to one quarter of the country.
This is the headline finding of this paper.  The cap is the
instrument by which redistribution of \emph{ownership} (not
\emph{flow}) becomes the policy outcome.

\paragraph{Productive-flow share.}  The productive-flow share
$\Phi$ rises from approximately 0.00 under status quo to 0.27
under the cap and 0.37 under cap-plus-tax.  $\Phi$ is the share of
per-tick monetary activity that goes to or is created by
productive-class agents (workers, makers, employers) rather than
extractive-class agents (rentiers, lenders, speculators).  Under
status quo, essentially all of the per-tick activity is the
rentier-to-lender deposit-pool loop and the capital-return
compounding within rentier stocks; the productive sector is
relatively inert in flow terms.  Under the cap, the structural
prevention of dynastic accumulation re-routes per-tick activity
through the productive classes.  The economy moves towards being
\emph{more productive in flow terms}, not just less unequal in
stock terms.  Cumulative value created over the 180-tick run rises
from 752 to 798 thousand model units, a 6\% increase.

\paragraph{Survival.}  Population survival at terminal tick rises
from 84\% under status quo to 91--92\% under the cap (with or
without the wealth tax).  The mechanism is the state-spending
package: revenue collected through inheritance excess or the
wealth tax is redistributed via UBI and state employment, which
catches working-class agents at the margin of subsistence
failure.  The cap delivers this survival improvement not by
making the cap itself life-saving but by funding the spending
package that is.

\paragraph{Revenue: the cap and the tax are complementary, not
competing.}  The wealth tax raises substantial state revenue:
the Zucman 2\% above \pounds 10M produces approximately
\pounds 56B per year at simulated UK scale, accumulating to
\pounds 5{,}165B in state wealth over the 180-tick run.  The
\pounds 5M cap raises far less revenue through the explicit
wealth-tax channel (zero, when paired with no wealth tax) but
accumulates state wealth of \pounds 484B through inheritance
excess at death.  Cap-plus-Zucman raises less than Zucman alone
(\pounds 0.8B/year and \pounds 335B cumulative) because the cap
prevents the buildup of wealth above the wealth-tax threshold that
would otherwise have been taxed; the two instruments compete for
the same upper-end tax base.

The honest reading is that the cap and the wealth tax are
\emph{instruments serving different goals}.  If the policy goal is
to fund the state at scale, the wealth tax is the instrument; the
cap reduces that revenue.  If the policy goal is to redistribute
ownership, the cap is the instrument; the wealth tax barely does
that.  Combining them produces the cleanest distributional
outcome at the cost of some wealth-tax revenue, and the policy
choice between cap-alone, tax-alone, and the combination is
fundamentally a choice between three different objectives.

%==========================================================================
\section{The right cap value: a sweep}
\label{sec:cap-value}
%==========================================================================

The choice of \pounds 5M as the cap value in Section~\ref{sec:comparison}
is the best-fit point from the joint calibration in our JEIC companion.
It does not follow that \pounds 5M is the correct cap value for any
particular policy goal.  We sweep the cap from \pounds 1M to no cap
(in nine log-spaced steps), both as a sole intervention and combined
with the Zucman 2\% tax.

\begin{figure}[t]
\centering
\includegraphics[width=\linewidth]{cap_sweep.png}
\caption{Cap-value sweep at the joint-calibrated UK baseline.  Four
metrics on the same x-axis (the cap value in pounds, log-spaced).
Solid line: cap alone.  Dashed line: cap plus Zucman 2\% above
\pounds 10M.}
\label{fig:capsweep}
\end{figure}

Two structural facts emerge.

\paragraph{The knee of the trade-off curve sits between \pounds 5M
and \pounds 10M.}  Above \pounds 10M the cap leaks: at \pounds 50M
the model produces 86 billionaires under the cap-alone arm and 6
under the cap-plus-Zucman arm; at \pounds 100M the leakage is
substantial (318 and 44 respectively); at \pounds 500M and above the
cap is essentially toothless.  Below \pounds 5M the cap begins to
catch middle-class succession and survival falls: from 83\% at
\pounds 5M to 80\% at \pounds 1M.  The model's stylised demography
is too coarse to put a confident lower bound on the cap, but the
direction of the survival cost as the cap is tightened is robust.

\paragraph{The wealth tax shifts the knee inward by a factor of
roughly two.}  At \pounds 5M cap plus Zucman 2\%, the top-1\% share
is 0.135, compared to 0.198 under cap alone.  At \pounds 10M cap
plus Zucman 2\%, the top-1\% share is 0.140 (compared to 0.222 under
cap alone) and the billionaire count is still zero.  The combination
allows a slightly higher cap to achieve the same top-share
compression as a tighter cap alone; this is the design space the
political debate could plausibly negotiate within.

\paragraph{Goal-dependent recommendations.}  The right cap value is
a function of which goal one is optimising for.

\begin{itemize}
\item \emph{Eliminate the billionaire class with minimal mortality
cost}: \pounds 5M to \pounds 10M, the knee of the trade-off curve.
This is where the model's qualitative finding is robust.

\item \emph{Match the UK Pareto slope}: \pounds 5M.  The Pareto
slope of the model under this configuration is $-0.45$, against
an empirical target of $-0.71$; the slope is reachable at this cap
value, the levels of top-1\% and top-0.1\% share remain below the
UK targets.

\item \emph{Minimum disruption to middle-class succession}:
\pounds 10M, with explicit carve-outs for productive business
assets and primary residences (Section~\ref{sec:implementation}).
At this cap the cap-alone arm leaves 6 billionaires and the
cap-plus-Zucman arm leaves 0; both numbers are an acceptable
trade for avoiding family-business succession damage.

\item \emph{Politically achievable in the UK alone before tax
flight kicks in}: \pounds 25M to \pounds 50M, paired with a hard
exit tax.  At \pounds 50M the cap-alone arm produces 86
billionaires; the cap-plus-Zucman arm produces 6.  This is the
compromise number, not the structurally optimal one.

\item \emph{G20-coordinated implementation}: \pounds 5M to
\pounds 10M.  Coordination removes the international flight
constraint and lets the model's structurally optimal range apply.
\end{itemize}

%==========================================================================
\section{Implementation}
\label{sec:implementation}
%==========================================================================

The model says nothing about how to implement the cap.  Implementation
is the harder political problem and we close it as honestly as we can
on the basis of the existing empirical and legal literature.

\subsection{The international flight problem}

The single-generation flight problem is the strongest unilateral-
implementation objection.  Under a UK-only inheritance cap, the
children of current UK billionaires have one generation (roughly
30 years on the parent's age at the law's passage) to establish
foreign tax residence.  The legal architecture is well-developed:
high-net-worth jurisdictions including Monaco, Singapore,
Switzerland, Dubai, and several others already function as
inheritance-tax havens, and the major UK private-client legal firms
have established expertise in restructuring assets to use them.

The empirical tax-flight literature
\citep{young2016millionaire, akcigit2016inventors, moretti2017taxing}
finds smaller flight responses than political opponents typically
forecast, but those studies measure mobility in response to marginal
rates on income or wealth, not in response to absolute caps on
inheritance.  An absolute cap is a structurally different instrument
because the trigger (death) is a one-shot event with predictable
timing in expectation, and because the response (becoming tax-
resident elsewhere) has to be settled before the trigger fires.
We expect the empirical elasticity of inheritance-cap response to
be higher than the existing studies' estimates of wealth-tax
response, but the literature does not currently let us put a number
on this.

Two design elements address the unilateral-implementation problem:

\paragraph{A residency clock.}  The cap should apply to any
individual who has been UK tax-resident for any of the previous 20
years, with the cap value reduced pro-rata for fewer years of
residency.  This makes the unilateral flight strategy expensive in
expected value because the lead time required to establish
non-coverage extends across a substantial fraction of a lifetime.
The architecture mirrors the UK's existing 17-out-of-20 ``deemed
domicile'' rule for inheritance tax but extends the coverage period
substantially.

\paragraph{An exit tax.}  The cap should crystallise at the point of
renouncing UK tax residence, not only at death.  This makes
renunciation itself the taxable event and removes the time-arbitrage
strategy.  The architecture mirrors the US exit-tax framework
(IRC Section 877A) which already taxes deemed-realisation at
expatriation; the difference is that the UK exit tax would apply
the cap rather than the standard estate-tax rate.

\paragraph{G20-coordinated implementation.}  The Zucman 2024
proposal pitched coordination at the G20 level as the
institutional innovation; the same logic applies, more strongly, to
an inheritance cap.  G20 jurisdictions hold approximately 80\%
of global wealth; a coordinated cap at the G20 level closes
substantially more of the flight surface than a UK-only cap.

\subsection{Trust law and asset definitions}

The historical effectiveness of inheritance taxation has been limited
by the use of trusts and other legal structures to separate beneficial
ownership from formal title.  A cap on \emph{inheritance} that did not
pierce trusts would be straightforwardly avoided.  Design elements:

\paragraph{Beneficial-ownership tracing.}  Assets held in trust at
the time of an individual's death (or expatriation) should be
attributed to that individual for cap purposes in proportion to
the individual's beneficial interest, with a residual rule that
attributes assets to the original settlor if beneficial interests
are otherwise undefined.  This is the same architecture as the
UK's existing trust-anti-avoidance regime for inheritance tax,
extended to absolute beneficial ownership.

\paragraph{Gift inter-vivos taper.}  Gifts made in the seven years
prior to death (or expatriation) should be counted toward the cap.
Existing UK IHT design includes a seven-year taper for this reason.
Whether the taper should be longer for cap purposes than for
existing IHT is a policy choice; the model is silent on this.

\paragraph{Life-insurance and pension carve-outs.}  Life insurance
proceeds and defined-benefit pensions are economically inheritance
in many cases and should be treated as such for cap purposes, with
explicit carve-outs only for fully-portable defined-contribution
pension pots up to a generous personal cap (perhaps the same
\pounds 5M).

\subsection{Asset-class carve-outs}

The case for asset-class carve-outs is empirical and political
rather than structural.  The model is too coarse to distinguish
productive business assets from financial wealth; it treats all
agent wealth as a single fungible state variable.  Real
implementation will need to make choices.

\paragraph{Productive-business assets.}  A sliding-scale carve-out
that protects the first \pounds X of productive-business assets and applies
the cap above this preserves family-business succession in the
small-and-medium-enterprise band.  The UK's existing Business
Relief regime is the architectural model; the design choice is
where to set X and what counts as a productive-business asset.
The empirical record on Business Relief is mixed; it has been used
extensively for tax-avoidance purposes alongside its intended
small-business-succession purpose.

\paragraph{Primary residences.}  A carve-out for the primary
residence up to a generous threshold (perhaps \pounds 2M for a UK
implementation, indexed to housing) preserves middle-class
intergenerational housing wealth.  The political economy of this
carve-out is straightforward: it is the carve-out that median UK
voters care about.

\paragraph{Charitable bequests.}  Existing UK charity-relief
regimes for inheritance already exist; the question for a cap
implementation is whether charity-funded family foundations
(the Gates Foundation architecture) should count toward the cap
or not.  This is a design choice that the model cannot adjudicate;
the political-economy case for treating large family foundations
as inheritance-equivalent is strong but the implementation is
fraught.

\subsection{Comparison to existing instruments}

The UK Inheritance Tax (40\% above \pounds 325k threshold) is the
nearest existing UK instrument and is, by international comparison,
unusually low both in threshold and rate.  The US Estate Tax (40\%
above \$13.6M per individual) has a higher threshold but operates
on the same rate-based architecture.  Both are rate-based, both
permit a large family-business carve-out, and both have generated
substantial avoidance infrastructure.

The proposed cap differs structurally.  It is an absolute ceiling,
not a marginal rate; it eliminates the ``one more pound'' arithmetic
that drives marginal-rate avoidance; and it forces the design
choice for the cap value into the political debate explicitly,
where in the existing rate-based system the design choice is
hidden in the marginal-rate schedule.

%==========================================================================
\section{Discussion}
\label{sec:discussion}
%==========================================================================

\subsection{Why is this not in any manifesto?}

The proposal of an absolute cap on inheritance is in no major
political manifesto in the UK as of the time of writing.  The
contemporary debate features the Zucman 2\% wealth tax in various
formulations, several different progressive wealth-tax schedules
(Warren, Sanders), various proposals to tighten the existing
Inheritance Tax (lower threshold, higher rate, fewer carve-outs),
and ongoing proposals to abolish IHT entirely.  None of these is
the absolute-cap proposal.

The political-economy explanation is, we suspect, two-fold.  First,
the cap is a structurally different instrument from the rate, and
the policy vocabulary of the contemporary debate is rate-based.
The cap is the instrument the analytics points to but the analytics
is in technical papers; the political debate operates at the level
of headline rates that voters can compare across proposals.

Second, the cap is the instrument that simultaneously raises the
ownership stake of the bottom half of the country and preserves
middle-class succession.  This combination does not fit either
ideological camp's natural framing.  Existing left proposals
attack the top with rates and accept that the bottom does not
gain ownership directly.  Existing right proposals defend the
middle with carve-outs and accept the top.  An absolute cap on
inheritance simultaneously moves wealth share to the bottom half
(from 0\% to 25\%) and leaves the middle largely untouched (a
\pounds 5M cap protects the 99th percentile of UK households
comfortably).  No major political programme is currently shaped to
make this combination a slogan.

We frame the case here in those terms, which we believe is the
correct framing for the contemporary political moment.  Calling
the proposal ``end the billionaire class'' is technically correct
but politically miscoded as vendetta against a particular
demographic.  Calling the proposal ``let the bottom half of the
country own something'' is the same instrument described by what
it does, which we believe lands better in the contemporary
discourse.

The instrument is also, in a sense, what economists would call
\emph{unfashionable in mechanism design}.  The rate-based optimal-
taxation tradition descends from
\citet{piketty2013theory, farhi2010progressive, saez2018social}
and operates in the language of marginal rates and elasticities.
An absolute cap maps awkwardly into this framework because the
marginal rate at the cap is undefined (or infinite).
\citet{atkinson2015inequality} is the rare recent reference that
proposes something cap-shaped, and even Atkinson framed it as a
universal capital endowment (a redistribution mechanism) rather
than as an absolute inheritance limit.

\subsection{Limits of the model}

The model is a stylised intuition pump, not a predictive instrument.
It has no spatial structure, no firm formation, no fiscal-crisis
dynamics for the state, toy speculator dynamics that do not produce
bubbles, and a single lender pool with no shadow-banking.  The
demographic assumptions (lifespan, age distribution at
initialisation) are crude.

The joint calibration matches the UK Pareto slope and the top-10\%
share at 80\% survival, which is a reasonable summary fit but is
not a fully validated calibration.  The remaining mismatches
(top-1\% and top-0.1\% share both below the UK targets) suggest
that the next calibration step should free the population
proportions (the rentier share at 3\% of the population is
plausibly too high for the strict billionaire class).  A joint
calibration that frees the population proportions and the lifespan
distribution to fit the wealth tail, the survival rate, and the
demographic age distribution simultaneously is the natural follow-up.

The cap-value sweep is at $N = 10{,}000$ with 4 seeds per cell.
Statistical robustness is adequate for the qualitative claims but
the confidence intervals on the specific cap values at the knee
are not tight.  A scale-up and seed-up pass is the natural next
robustness step.

The model is also entirely closed-economy.  The international
flight question is treated qualitatively here on the basis of the
empirical literature and the implementation-architecture arguments;
the model does not produce a quantitative estimate of the cap
elasticity to international tax differentials.  An extension of the
model to include a foreign residence option for agents above a
wealth threshold is a natural next step.

%==========================================================================
\section{Conclusion}
%==========================================================================

We have presented a calibrated agent-based simulation of a
stylised UK-like economy and used it to compare five fiscal-policy
interventions on three structural outcomes: ownership distribution,
productive-flow share, and survival.  The headline finding is that
an absolute \pounds 5M cap on inheritance raises the wealth share
held by the bottom half of the population from 0.0\% under status
quo to 19.2\% under the cap alone and 25.3\% under cap-plus-tax.
The cap is the only instrument among the five we tested that
changes who owns the country.

The cap also raises the productive-flow share $\Phi$ from
approximately zero under status quo to 0.27 under the cap and 0.37
under cap-plus-tax, increases cumulative value created by about
6\%, and improves population survival from 84\% to 91--92\%.  The
disappearance of the billionaire class is the visible side effect
of these structural changes, not the policy goal.

A cap-value sweep places the knee of the trade-off curve between
\pounds 5M and \pounds 10M.  The right value within this range is
a political choice that depends on whether the priority is the
redistribution of ownership (lower cap) or the protection of
middle-class business succession (higher cap, with carve-outs).

The argument extends but does not displace the Zucman framing.  A
2\% G20-coordinated wealth tax is an excellent revenue instrument
and the case for international coordination is strong.  Our
amendment is that the G20 coordination should be on the cap as
well as on the rate, and that the cap is a categorically
different policy lever: the tax raises revenue at scale, the cap
redistributes ownership.  The two are complementary not competing.
The cap has the deeper philosophical lineage
\citep{mill1848principles, carnegie1889wealth, rignano1924bequest,
atkinson2015inequality} and the wealth tax has the contemporary
political momentum.

We close with the positive claim:
\emph{a \pounds 5 million to \pounds 10 million absolute cap on
inheritance, internationally coordinated at the G20 level, with
residency-clock and exit-tax provisions and explicit carve-outs
for primary residences and small productive-business assets, is
the single most effective structural fiscal instrument the model
identifies for raising the ownership stake of the bottom half of
the population from zero to roughly one quarter of the country's
wealth, increasing the share of monetary activity that is
productive rather than extractive, and improving population
survival.}  The end of the billionaire class is the visible side
effect.  The proposal is about ownership.

\bibliographystyle{plainnat}
\bibliography{references}

\appendix
\section{Cap-value sweep data}
\label{app:sweep}

\begin{table}[h]
\centering
\small
\caption{Full results of the cap-value sweep at the joint-calibrated
UK baseline.  $N = 10{,}000$, 4 seeds, 180 ticks.  Solid line in
Figure~\ref{fig:capsweep}.}
\begin{tabular}{lrrrrrrr}
\toprule
arm & cap & alive \% & $\geq$\pounds 1B & $\geq$\pounds 100M
& $\geq$\pounds 10M & top-1\% & Pareto slope \\
\midrule
cap alone & \pounds 1M   & 80\% &   0 &  27 &   996 & 0.189 & $-0.698$ \\
cap alone & \pounds 5M   & 83\% &   0 & 208 & 1{,}424 & 0.198 & $-0.455$ \\
cap alone & \pounds 10M  & 83\% &   6 & 403 & 4{,}016 & 0.222 & $-0.454$ \\
cap alone & \pounds 50M  & 83\% &  86 & 721 & 4{,}084 & 0.323 & $-0.443$ \\
cap alone & \pounds 100M & 83\% & 318 & 840 & 4{,}084 & 0.352 & $-0.442$ \\
cap alone & \pounds 500M & 83\% & 700 & 848 & 4{,}084 & 0.391 & $-0.440$ \\
cap alone & \pounds 1B   & 83\% & 803 & 848 & 4{,}084 & 0.405 & $-0.440$ \\
cap alone & no cap        & 83\% & 804 & 848 & 4{,}084 & 0.327 & $-0.260$ \\
\midrule
cap + Zucman 2\% & \pounds 1M   & 80\% &   0 &  16 &   980 & 0.146 & $-0.540$ \\
cap + Zucman 2\% & \pounds 5M   & 83\% &   0 &  62 & 1{,}414 & 0.135 & $-0.333$ \\
cap + Zucman 2\% & \pounds 10M  & 83\% &   0 &  98 & 3{,}998 & 0.140 & $-0.321$ \\
cap + Zucman 2\% & \pounds 50M  & 83\% &   6 & 593 & 4{,}057 & 0.230 & $-0.296$ \\
cap + Zucman 2\% & \pounds 100M & 83\% &  44 & 804 & 4{,}057 & 0.272 & $-0.291$ \\
cap + Zucman 2\% & \pounds 500M & 83\% & 497 & 804 & 4{,}057 & 0.335 & $-0.282$ \\
cap + Zucman 2\% & \pounds 1B   & 83\% & 800 & 804 & 4{,}057 & 0.353 & $-0.263$ \\
cap + Zucman 2\% & no cap        & 83\% & 800 & 804 & 4{,}057 & 0.293 & $-0.226$ \\
\bottomrule
\end{tabular}
\end{table}

\end{document}
